% !TeX root = ../PIFB2.tex
\section{Assumptions, Results and Open Obligations}\label{sec:epistemic_status}

The framework has a small set of primitive choices: a continuous base and local domains; a gauge action on a statistical sample space; normalized priors, record kernels and interaction channels; a recognition family; and, when dynamics are discussed, a metric and mobility or a stochastic generator. Coarse-graining additionally requires a retained channel. An episode-history law for latent links is an additional normalized model with a declared increment penalty. Its Euclidean stationarity does not supply a real-time kinetic action, Lorentz soldering, thermodynamic interpretation or physical ontology.

Most mathematical results are conditional identities or constructions. An evidence bound follows from a normalized joint, source attention follows from an explicit label model, and coarse actions follow from integration. None of these identities selects the joint, the source alphabet or the retention by itself. The distinction lets the theory use exact results without claiming that the variational principle derives every modeling ingredient.

\begin{table}[ht]
\centering\small
\begin{tabularx}{\textwidth}{@{}p{0.27\textwidth}XX@{}}
\toprule
Claim & Established scope & Remaining obligation\\
\midrule
Selective communication & Endpoint covariance, rank bounds, normalized noisy channels. & Learned interface identifiability and predictive meaning.\\
Collective VFE and attention & Exact posterior identity and conditional or restricted source rows. & A generative model for a target system.\\
Latent interaction links & Normalized generator prior, Fisher recognition flow and controlled Euclidean Proca-type history equation. & Coarse link-law closure, rate separation and physical identification.\\
One statistical bundle & Conditional state--model decomposition and exact VFE accounting. & Emergent model coordinates and their slow kinetics.\\
Persistent parents & Exact pushforwards, descent conditions and memory equations. & Finite predictive closure for the actual evolving network.\\
Formation & Fixed-temperature four-cycle pairing selection in the companion model. & General selective-interface and recursive formation.\\
Information clocks & Fisher length, dissipation identity and fixed-channel contraction. & Physical time calibration and causal order.\\
Lorentz structure & Hermitian determinant form and the $SL(2,\mathbb C)$ double cover. & Dynamic channel selection, soldering and physical propagation.\\
Mass and persistence & Response, stiffness, memory and metastable lifetimes in stated models. & A derived inertial mass law and physical units.\\
\bottomrule
\end{tabularx}
\caption{The distinction between proved conditional mathematics and the research program.}\label{tab:claim_status}
\end{table}

The technical appendices sometimes use the short status codes P for a proved conditional claim, D for an explicit derivation, S for a structural correspondence, C for a conjecture or open obligation, and E for empirical evidence. These codes classify statements rather than certify a whole section. No empirical realization of the physical program is reported here.

\subsection{Scope and limitations}\label{sec:scope_limitations}
Formation and inheritance are analyzed at one fixed base point. The smooth base remains the ambient setting; conditional base geometry appears only when a derivative or soldering is explicitly supplied. Gaussian families, fixed links, compact gauge reductions, prescribed partitions and imposed rate separation are named special cases. They are not assumptions of every result.

A normalizer matters whenever parameters are varied. A posterior from one inference problem cannot be silently reused as a likelihood in another. A small marginal divergence need not control derivatives, force errors or path laws. A successful first parent does not establish the next formation episode. These are the principal sources of overextension the revised construction avoids.

\subsection{Relation to the mathematical counterpart}\label{sec:companion_theory}
The companion \emph{Gauge-Covariant Variational Free Energy and Exact Coarse-Graining} develops the measure-theoretic and geometric statements in detail~\cite{DennisGaugeVFE2026}. Its Gaussian supplement provides exact finite-dimensional examples~\cite{DennisGaugeVFESuppG2026}; its base supplement supplies the separate assumptions for transport, infinite networks and spatial realizations~\cite{DennisGaugeVFESuppB2026}. This monograph gives the physical motivation, worked reasoning and interpretive questions. It neither treats the companion's prototypes as the only possible agents nor treats an unresolved prototype as a general impossibility theorem.

